\section{Discussion}
The results of the evaluation reveal a precise structural tradeoff inherent in the Grounded Information Generation (GIG) approach: moving from unstructured similarity-based retrieval to strict ontology-grounded reasoning successfully suppresses hallucination and enforces Open World abstention, but it incurs a severe penalty in answer recall due to cumulative extraction and resolution failures. This section critically examines these findings, contextualizing the performance gap between ideal (ontology-grounded) and realistic (source-grounded) conditions, analyzing the failure modes within the pipeline, and assessing the validity of the RAG baseline comparison.

\subsection{The Cost of Grounding: Ideal vs. Realistic Performance}
The most striking finding is the disparity between the system's performance on the ontology-grounded dataset (Table \ref{tab: ontology-grounded E2E}) and the source-grounded dataset (Table \ref{tab: source-grounded GIG}). When evaluated on questions guaranteed to correspond to explicitly extracted ontology facts, the GIG system is highly robust. It produces no unacceptable groundings (0.0\%), and correctly abstains when facts are missing. Once the entity is resolved, answer generation succeeds almost universally (93.8\% acceptable grounding).

However, when evaluated against realistic questions derived directly from the source text, performance degrades sharply. The system yields only 24.4\% correct answers, dominated instead by a 48.9\% incorrect abstention rate. This drop demonstrates that while the \textit{inference} subsystem is sound, the \textit{knowledge-base construction} subsystem fails to capture or correctly consolidate the full semantic breadth of the source text. If a fact is missed during the offline chunk-wise extraction phase, or if an entity is merged incorrectly during the aggressive instance-resolution cleanup, that information is permanently lost to the inference engine. 

Consequently, the GIG system behaves exactly as designed under the Open World Assumption (OWA): when it cannot find a strict, schema-compliant path between resolved entities, it abstains. The high incorrect abstention rate is not a symptom of reasoning failure, but rather the intended failure mode of a strictly constrained system operating over an incomplete knowledge base.

\subsection{Pipeline Brittleness and Cascading Failures}
The high rate of incorrect abstentions is exacerbated by the decomposition-first architecture. While single-task decomposition improves interpretability—allowing us to trace exactly which component failed—it creates a brittle, cascading pipeline. For the system to answer a question correctly, it must succeed at every sequential step: question-type classification, entity extraction, relation mapping, object extraction, entity resolution, and final answer synthesis. 

The component-level evaluation (Table \ref{tab: Interpret Input}) highlights this vulnerability. While early stages like atomic-to-graph translation succeed at a 95.9\% rate, downstream object extraction strict-matching falls to 60.1\%. Furthermore, the deterministic relation mapping limits the system's expressiveness; if the natural language question implies a relation that does not map perfectly to the single allowed predicate for that question type, the retrieval stage naturally yields an empty context. 

Unlike a monolithic LLM, which can use latent semantic intuition to "bridge the gap" when input is slightly ambiguous, the GIG system's strict I/O schemas force a hard failure. This confirms a fundamental challenge in neuro-symbolic design: integrating symbolic rigidity with neural models effectively stops hallucinations, but the rigidity itself becomes the primary bottleneck to system flexibility.

\subsection{RAG Baseline Comparison and Fair Evaluation}
Comparing the GIG system directly to the RAG baseline (Figure \ref{fig:source-grounded-comparison}) clarifies the broader tradeoff between the two paradigms. The RAG baseline achieves a higher correct answer rate (45.5\% vs. 24.4\%) but suffers from a severe incorrect-false hallucination rate (36.7\% vs. 16.3\%) and fails completely at correct abstention (0.0\% vs. 75.0\%).

It must be acknowledged that this comparison contains a structural confound. The RAG baseline was restricted to a $k=3$ retrieved chunk limit. For aggregation questions (e.g., Membership or Capability) where the answer spans much of the document, the RAG system is starved of context, whereas the GIG system aggregated this information offline into the global ontology. Conversely, RAG possesses a distinct advantage in semantic recall: it retrieves raw text and uses the LLM to interpret it on the fly, skipping the lossy, loss-inducing intermediate step of ontology extraction entirely.

Despite these confounds, the comparison validates the thesis hypothesis. The current "try-your-best" generative paradigm (RAG) prioritizes answering over accuracy, resulting in frequent and confidently stated false information. GIG effectively eliminates these structural hallucinations and honors the limits of its knowledge, proving that strict schema grounding is a viable strategy for high-precision, low-hallucination environments—provided the upstream ontology extraction can be made sufficiently exhaustive.

\subsection{Limitations and Future Work}
The primary limitation of the current instantiation is the lossiness of the offline knowledge-base construction. Future iterations of the GIG approach must address the recall bottleneck at the extraction phase. Potential solutions include iterative, multi-pass ontology extraction over overlapping chunks, or utilizing a more expressive set of base properties to capture subtle semantic nuances that the current 14-class, 25-axiom DUL subset forces into overly generic bins.

Furthermore, the "Atomic Question Assumption" currently limits the system to single-hop reasoning. While this was necessary to establish a baseline of deterministic reliability, a fully mature system must implement an upstream agent capable of decomposing complex, multi-hop natural language queries into a sequence of atomic queries, feeding them through the GIG pipeline, and aggregating the results.
